\section{Anforderungen und methodisches Vorgehen}\label{sec:Anforderungen}

\subsection{Anwendungsszenarien}

Die Anwendungsszenarien ergeben sich aus dem Ziel, Robotermissionen unmittelbar mit einem digitalen
Gebäudemodell zu verknüpfen. Ausgangspunkt ist daher stets ein IFC-Gebäudemodell, das im
browserbasierten Editor geladen und räumlich betrachtet wird. Die dreidimensionale Darstellung dient
dabei nicht nur der Visualisierung, sondern bildet den räumlichen Bezugspunkt für die Beschreibung
der einzelnen Aufgaben einer Mission.

Ein grundlegendes Szenario besteht in der Auswahl eines konkreten Gebäudeelements als Bezug einer
Roboteraufgabe. Soll ein Roboter beispielsweise eine Tür öffnen, zu einem technischen Gebäudeelement
navigieren oder einen Schalter bedienen, wählt der Benutzer das betreffende Element direkt in der
Gebäudedarstellung aus und ordnet ihm eine entsprechende semantische Aktion zu. Der Aufgabenbezug
bleibt dadurch nicht auf eine abstrakte Positionsangabe beschränkt, sondern ist mit einem eindeutig
identifizierbaren Bestandteil des Gebäudemodells verbunden.

Darauf aufbauend können mehrere Aufgaben zu einer geordneten Robotermission zusammengestellt werden.
Ein typischer Ablauf kann beispielsweise darin bestehen, zunächst zu einer Tür zu navigieren, diese
zu öffnen, den Durchgang zu passieren und die Tür anschließend wieder zu schließen. In gleicher
Weise lassen sich weitere Gebäudeelemente in den Missionsablauf einbeziehen. Die einzelnen Aufgaben
beschreiben dabei fachlich, was im Zusammenhang mit einem bestimmten Gebäudeelement erfolgen
soll.

Ein weiteres Anwendungsszenario betrifft die Weiterbearbeitung bereits annotierter Gebäudemodelle.
Wird eine zuvor bearbeitete IFC-Datei erneut geladen, können die darin enthaltenen
Missionsinformationen rekonstruiert und im Editor wieder dargestellt werden. Der Benutzer kann
bestehende Aufgaben ändern, ergänzen, entfernen oder neu anordnen und die überarbeitete Mission
anschließend erneut als annotierte IFC-Datei exportieren. Die IFC-Datei dient damit sowohl als
Gebäudemodell als auch als fachlicher Träger der zugehörigen Missionsannotation.

\subsection{Funktionale Anforderungen}
Aus den zuvor beschriebenen Anwendungsszenarien ergeben sich die funktionalen Anforderungen an den
IFC-Editor. Sie beschreiben die aus Nutzersicht erforderlichen Systemfunktionen für die
Erstellung, Bearbeitung und persistente Verknüpfung von Robotermissionen mit einem Gebäudemodell.

\textbf{Laden und Darstellen von Gebäudemodellen}\newline
Das System muss IFC-Gebäudemodelle laden und räumlich darstellen können. Die Gebäudedarstellung
bildet die Grundlage dafür, Bauteile und andere geeignete IFC-Objekte als Bezug einer Roboteraufgabe
zu identifizieren und für die weitere Annotation auszuwählen.

\textbf{Gezielte Auswahl von Gebäudeelementen}\newline
Gebäudeelemente müssen zuverlässig für die Annotation auswählbar sein. Da sich in der
dreidimensionalen Darstellung mehrere Elemente überlagern können, muss das System auch bei mehreren
möglichen Treffern eine gezielte Auswahl des tatsächlich gewünschten Gebäudeelements ermöglichen.

\textbf{Verwaltung von Robotermissionen}\newline
Robotermissionen müssen erstellt, bearbeitet und gelöscht werden können. Eine Mission muss aus
mehreren Einzelaufgaben aufgebaut werden können, deren Reihenfolge festgelegt und nachträglich verändert
werden kann. Dadurch lassen sich mehrstufige Abläufe innerhalb des Gebäudes als zusammenhängende
Mission beschreiben.

\textbf{Beschreibung semantischer Roboteraufgaben}\newline
Für die einzelnen Einzelaufgaben müssen die für ihre fachliche Bedeutung erforderlichen Informationen erfasst
und bearbeitet werden können. Dazu gehört insbesondere die Zuordnung einer semantischen
Roboteraktion, durch die beschrieben wird, welche Aufgabe im jeweiligen Missionsschritt ausgeführt
werden soll.

\textbf{Verknüpfung von Einzelaufgaben mit dem Gebäudemodell}\newline
Einzelaufgaben müssen mit geeigneten IFC-Gebäudeelementen verknüpft werden können. Der Objektbezug stellt
sicher, dass eine Aufgabe nicht nur abstrakt beschrieben wird, sondern eindeutig auf den
betreffenden Bestandteil des Gebäudes verweist.

\textbf{Validierung von Missionen und Einzelaufgaben}\newline
Das System muss die erfassten Missions- und Aufgabeninformationen auf fachliche Konsistenz prüfen
können. Hierzu gehören insbesondere erforderliche Angaben, gültige Objektbezüge und eine
widerspruchsfreie Anordnung beziehungsweise Abhängigkeit der einzelnen Tasks. Erkannte Probleme
müssen vor einer weiteren Verarbeitung nachvollziehbar kenntlich gemacht werden.

\textbf{Export als annotierte IFC-Datei}\newline
Eine erstellte oder bearbeitete Robotermission muss einschließlich ihrer Einzelaufgaben, Aktionsinformationen
und Objektzuordnungen in einer annotierten IFC-Datei gespeichert beziehungsweise exportiert werden
können. Die Missionsbeschreibung bleibt dadurch mit dem zugrunde liegenden Gebäudemodell in einer
gemeinsamen Austauschdatei verbunden.

\textbf{Rekonstruktion und erneute Bearbeitung annotierter Missionen}\newline
Bereits annotierte IFC-Dateien müssen erneut geladen werden können, wobei die darin enthaltenen
Robotermissionen rekonstruiert werden. Diese Missionen müssen anschließend weiterbearbeitet und
erneut exportiert werden können. Wiederholte Bearbeitungs- und Exportzyklen dürfen dabei nicht zu
einer unkontrollierten Vervielfachung der vom System erzeugten Missionsannotation führen. 


\subsection{Nichtfunktionale Anforderungen}

Neben den funktionalen Anforderungen bestimmen qualitative Randbedingungen, unter welchen
Voraussetzungen der Viewer und Editor zuverlässig und praktikabel eingesetzt werden kann. Im

\textbf{Browserbasierte Nutzbarkeit  }\newline
Der Viewer und Editor soll in einer modernen Browserumgebung ausführbar sein und keine Installation
einer spezialisierten Desktop-BIM-Anwendung voraussetzen.

\textbf{Darstellungsperformance  }\newline
Auch umfangreichere IFC-Gebäudemodelle sollen so verarbeitet und dargestellt werden können, dass
Navigation, Objektauswahl und Annotation nutzbar bleiben.

\textbf{Stabile Referenzierung  }\newline
Die fachliche Verbindung zwischen RobotTasks und den zugehörigen Gebäudeelementen soll über
Bearbeitungs-, Speicher- und Exportzyklen erhalten bleiben. 

\textbf{Datenintegrität und Roundtrip-Stabilität  }\newline
Das erneute Laden, Bearbeiten und Exportieren einer annotierten IFC-Datei darf die enthaltene
Missionssemantik nicht unbeabsichtigt verändern oder Missionsinformationen durch wiederholte
Exportvorgänge duplizieren. Gleichzeitig sollen Inhalte des Gebäudemodells, die nicht durch die
Missionsannotation verwaltet werden, unbeeinträchtigt bleiben.

\textbf{Erweiterbarkeit und Trennung der Verantwortlichkeiten  }\newline
Die fachliche Beschreibung einer Robotermission soll ausreichend von Darstellung, Persistenz und
späterer Roboterausführung getrennt sein. Dadurch sollen weitere Task-Arten sowie nachgelagerte
Funktionen wie roboterspezifische Exporte, Navigationsziele oder eine serverseitige Verarbeitung
ergänzt werden können, ohne die grundlegende fachliche Modellierung ersetzen zu müssen.

\textbf{Nachvollziehbarkeit und Robustheit  }\newline
Unvollständige oder widersprüchliche Missionen sollen erkennbar sein und zu nachvollziehbaren
Validierungshinweisen führen. Zustände, bei denen beispielsweise erforderliche Objektbezüge oder
zulässige Abhängigkeiten nicht eindeutig bestimmt werden können, sollen nicht unbemerkt in einen
Export übernommen werden.

\textbf{Benutzbarkeit der Gebäudezuordnung  }\newline
Die Zuordnung einer Roboteraufgabe zu einem konkreten Gebäudeelement soll auch bei komplexer oder
überlagernder IFC-Geometrie verständlich und kontrollierbar bleiben. Für den Benutzer muss
insbesondere nachvollziehbar sein, welches Bauteil als fachlicher Bezug einer Annotation ausgewählt
wurde.

